\chapter{Monte Carlo Statistical Analysis}
\label{ch:monte-carlo}

\section{Variation Model}
\label{sec:mc-model}

Following \cref{eq:mc-delay}, each Monte Carlo sample draws one global
random variable $Z_g$, shared by every gate in that sample, and one local
random variable $Z_i$ per gate, drawn independently. Because nominal gate
delay (\cref{eq:drise,eq:dfall}) is always strictly positive for any
physically valid cell, a sample is discarded and re-drawn only when the
\emph{multiplicative} variation factor $1 + \sigma_g Z_g + \sigma_l Z_i$
itself would be non-positive for some gate --- a condition that can be
checked without reference to any particular circuit's nominal delay values.

\section{Paired Pre-/Post-Optimization Comparison}
\label{sec:mc-pairing}

A specific correctness property is enforced by construction: the same set
of $M$ valid variation samples --- the same sequence of $(Z_g^{(k)},
Z_i^{(k)})$ draws, including the same accept/reject decisions from the
resampling rule --- is generated exactly once per run and then applied
unchanged to \emph{both} the pre-optimization circuit and the
post-optimization circuit. This is possible precisely because the
resampling condition above depends only on the variation factor, not on
any circuit-specific nominal delay, so the same sample set is guaranteed
valid for both circuits regardless of how their sizing differs. The result
is a genuinely paired comparison: sample $k$ represents the \emph{same}
hypothetical manufactured die under both sizing assignments, which is a
meaningfully stronger guarantee than independently reseeding a random
number generator for each of the two runs, since two independently-seeded
runs are only guaranteed to draw identically if they consume randomness in
exactly the same order and quantity, which is not obviously true once
resampling is involved and the two circuits' nominal delays differ.

\section{Reported Metrics}
\label{sec:mc-metrics}

For each of the two paired runs, full STA (\cref{ch:sta}) is re-executed
once per sample against that sample's perturbed delay set, and the
following are reported: the mean, standard deviation, minimum, and maximum
of circuit delay, WNS, and TNS across all samples; the empirical violation
probability and timing yield (\cref{eq:mc-yield}); the modal (most
frequently observed) critical path and the fraction of samples on which it
occurs, as a measure of critical-path stability under variation; and, per
gate, the empirical probability that it appears on the sample's single most
critical path, which identifies gates whose criticality is stable versus
gates that only become critical under specific variation draws.

\section{Interpretation}
\label{sec:mc-interpretation}

Comparing the pre- and post-optimization statistical summaries answers a
question that a single, nominal-corner STA run cannot: whether the
gate-sizing changes made by the optimizer (\cref{ch:optimization}) improved
robustness to manufacturing variation, not merely the nominal-case WNS/TNS.
A design can be nominally compliant yet statistically fragile (a high
violation probability under realistic variation) if its post-optimization
slack margin is thin; \cref{ch:results} reports this comparison for the
supplied Monte Carlo-enabled benchmark circuit and discusses whether the
optimizer's nominal-case improvement translated into a statistically
meaningful yield improvement.
